%%=============================================================================
%% Voorwoord
%%=============================================================================

\chapter*{\IfLanguageName{dutch}{Woord vooraf}{Preface}}%
\label{ch:voorwoord}

%% TODO:
%% Het voorwoord is het enige deel van de bachelorproef waar je vanuit je
%% eigen standpunt (``ik-vorm'') mag schrijven. Je kan hier bv. motiveren
%% waarom jij het onderwerp wil bespreken.
%% Vergeet ook niet te bedanken wie je geholpen/gesteund/... heeft

Voor u ligt mijn bachelorproef, geschreven tijdens mijn opleiding Toegepaste Informatica aan HOGENT. In deze bachelorproef onderzoek ik hoe betrouwbaar large language models zijn bij het genereren van Python code voor data analyse en ETL taken. Ik heb bewust voor dit onderwerp gekozen omdat generatieve AI vandaag een steeds grotere rol speelt binnen softwareontwikkeling en data engineering, en omdat ik wilde onderzoeken in welke mate deze modellen werkelijk bruikbare, correcte en veilige code kunnen opleveren in realistische scenario's.

Tijdens mijn opleiding merkte ik dat tools zoals ChatGPT en andere coding assistants steeds vaker gebruikt worden om programmeerproblemen sneller op te lossen. Hoewel deze tools vaak indrukwekkende resultaten geven, stelde ik mij ook de vraag hoe correct die gegenereerde code werkelijk is, zeker wanneer het gaat over taken zoals datatransformatie, validatie, performantie en beveiliging. Met deze bachelorproef wilde ik daarom niet alleen modellen met elkaar vergelijken, maar ook beter begrijpen waar hun sterktes en beperkingen liggen binnen een data engineering context.

Ik wil in het bijzonder graag mijn promotor meneer Decorte en mijn copromotor meneer Vanhaute, bedanken voor hun begeleiding. ondersteuning en feedback tijdens het uitwerken van deze bachelorproef. Hun opmerkingen en suggesties hebben mij geholpen om mijn onderzoek beter af te stemmen, mijn methodologie verder uit te werken en mijn resultaten beter te begrijpen. Dankzij hun begeleiding kon ik deze bachelorproef inhoudelijk sterker en concreter maken.

Ik hoop dat deze bachelorproef een nuttige bijdrage vormt voor studenten en docenten die op een kritische manier met generatieve AI en codegeneratie willen omgaan.